% ============================================================
% Title supplied by appendix wrapper.
\label{sec:advanced-operators}
% ============================================================

\begin{quote}\small
This appendix provides expanded definitions, proofs, and examples for selected operator tools.
\end{quote}

This section introduces four reusable operator tools that package common derivations in Coherence Calculus:
horizon flow (discrete update across scales), spectral selection (spectral selector), trace aggregation (coherence trace), and canonical leakage direction (aliasing operator).

% ------------------------------------------------------------
\subsection{Operator Reference}
\label{subsec:operator-atlas}
% ------------------------------------------------------------

This appendix provides expanded definitions, proofs, and examples for the operators introduced in the main text.

% ------------------------------------------------------------
\subsection{Horizon Flow Operator (Discrete Update)}
\label{subsec:horizon-flow}
% ------------------------------------------------------------

The horizon flow operator is a \emph{discrete update} that measures how an observed operator changes across adjacent coherence scales.
(Note: this is a finite difference, not the tower calculus derivative $\partial = [\CCIndexOp, \cdot]$.)

\begin{definition}[Horizon-restricted operator]
\label{def:horizon-restricted}\lean{CoherenceCalculus.AppendixOperators.horizonRestrict}
For an operator $A:\cH \to \cH$ and horizon projection $\CCPi{h}$, define the \emph{horizon-restricted operator}
\[
\CCObs{A}{h} := \CCRestr{A}{h}.
\]
This is the operator ``seen'' by an observer at horizon level $h$.
\end{definition}

\begin{definition}[Horizon flow operator]
\label{def:horizon-flow}\lean{CoherenceCalculus.AppendixOperators.horizonFlow}
The \emph{horizon flow operator} is the discrete finite difference
\[
\CCFlow{A}{h} := \CCObs{A}{h+1} - \CCObs{A}{h} = \CCRestr{A}{h+1} - \CCRestr{A}{h}.
\]
This measures the change in the observed operator when the horizon is refined from $h$ to $h+1$.
\end{definition}

\begin{proposition}[Properties of horizon flow]
\label{prop:horizon-flow-linear}\lean{CoherenceCalculus.AppendixOperators.horizonFlow_add, CoherenceCalculus.AppendixOperators.horizonFlow_id_eq_increment, CoherenceCalculus.AppendixOperators.horizonFlow_proj_zero}
The horizon flow operator satisfies:
\begin{enumerate}[label=(\roman*)]
\item \textbf{Linearity:} $\CCFlow{(\alpha A + \beta B)}{h} = \alpha \CCFlow{A}{h} + \beta \CCFlow{B}{h}$ for scalars $\alpha, \beta$.
\item \textbf{Identity:} $\CCFlow{I}{h} = \CCPi{h+1} - \CCPi{h} = \CCDelta{h+1}$.
\item \textbf{Horizon projector:} $\CCFlow{\CCPi{h}}{h} = 0$.
\end{enumerate}
\end{proposition}

\begin{proof}
(i) Direct computation using linearity of the projections $\CCPi{h}$ and $\CCPi{h+1}$.

(ii) For the identity operator $I$, we have $\CCObs{I}{h} = \CCPi{h}$ and $\CCObs{I}{h+1} = \CCPi{h+1}$.
Thus $\CCFlow{I}{h} = \CCPi{h+1} - \CCPi{h}$.

(iii) For the horizon projector $\CCPi{h}$ viewed as an operator: $\CCObs{(\CCPi{h})}{h} = \CCPi{h}$.
For the refined restriction, $\CCObs{(\CCPi{h})}{h+1} = \CCPi{h}$ by transitivity.
Thus $\CCFlow{\CCPi{h}}{h} = \CCPi{h} - \CCPi{h} = 0$.
\end{proof}

\begin{quote}\small
\textbf{Intuition.}
The horizon flow $\CCFlow{A}{h}$ is a discrete finite-difference operator, not
the tower calculus derivative.
It measures how the observed operator changes when the horizon is refined by one level.
No external parameter is involved; only the internal structure of the horizon
tower enters.
(For the true calculus derivative with Leibniz rule and fundamental theorem, see the tower derivative $\CCDer{\cdot} = [\CCIndexOp, \cdot]$ in Section~\ref{subsec:tower-derivative}.)
\end{quote}

\begin{remark}[Continuous scale derivative]
\lean{CoherenceCalculus.AppendixOperators.continuousScaleDerivative}
If the horizon tower is indexed continuously by $\ell \in \bbR$ (as in Section~\ref{sec:log-periodic}), one may define the \emph{scale derivative}
\[
\partial_\ell A_\ell := \lim_{\epsilon \to 0} \frac{A_{\ell+\epsilon} - A_\ell}{\epsilon}
\]
in the strong operator topology, when this limit exists.
This is a derivative with respect to \emph{scale}, not with respect to any external parameter.
\end{remark}

% ------------------------------------------------------------
\subsection{Spectral Selector (Projection onto Discrete Spectrum)}
\label{subsec:quantization}
% ------------------------------------------------------------

The spectral selector projects onto eigenstates corresponding to a chosen discrete subset of the spectrum.

\begin{definition}[Spectral selector]
\label{def:quantization-projector}\lean{CoherenceCalculus.AppendixOperators.spectralSelector}
Let $A:\cH \to \cH$ be a self-adjoint operator with spectral measure $E_A$.
For a discrete subset $S \subseteq \bbR$, define the \emph{spectral selector} (also called \emph{quantization projector} in physics contexts)
\[
Q_S(A) := E_A(S) = \mathbf{1}_S(A),
\]
where $\mathbf{1}_S$ is the indicator function of $S$, applied via functional calculus.
\end{definition}

\begin{remark}[Finite-dimensional case]
\lean{CoherenceCalculus.AppendixOperators.finite_quantization_example}
In finite dimensions, if $A$ has eigenvalues $\lambda_1, \ldots, \lambda_n$ with orthonormal eigenvectors $v_1, \ldots, v_n$, then
\[
Q_S(A) = \sum_{j:\, \lambda_j \in S} v_j v_j^*.
\]
This is the orthogonal projection onto the span of eigenvectors with eigenvalues in $S$.
\end{remark}

\begin{proposition}[Properties of spectral selector]
\label{prop:quantization-properties}\lean{CoherenceCalculus.AppendixOperators.spectralSelector_idem, CoherenceCalculus.AppendixOperators.spectralSelector_meet, CoherenceCalculus.AppendixOperators.spectralSelector_disjoint_join, CoherenceCalculus.AppendixOperators.spectralSelector_whole, CoherenceCalculus.AppendixOperators.spectralSelector_empty}
For self-adjoint $A$ and discrete $S, T \subseteq \bbR$:
\begin{enumerate}[label=(\roman*)]
\item $Q_S(A)$ is an orthogonal projection: $Q_S(A)^2 = Q_S(A) = Q_S(A)^*$.
\item $Q_S(A) Q_T(A) = Q_{S \cap T}(A)$.
\item $Q_S(A) + Q_T(A) = Q_{S \cup T}(A)$ if $S \cap T = \emptyset$.
\item $Q_\bbR(A) = I$ and $Q_\emptyset(A) = 0$.
\end{enumerate}
\end{proposition}

\begin{proof}
These are standard properties of spectral projections from the spectral theorem for self-adjoint operators.
\end{proof}

\begin{example}[Finite-dimensional spectral selection]
\lean{CoherenceCalculus.AppendixOperators.finite_quantization_example}
\label{ex:finite-quantization}
Let $\cH = \bbR^4$ and $A = \mathrm{diag}(0, 1, 2, 3)$ with standard basis $e_1, e_2, e_3, e_4$.
For $S = \{0, 2\}$:
\[
Q_S(A) = e_1 e_1^* + e_3 e_3^* = \mathrm{diag}(1, 0, 1, 0).
\]
This projects onto the subspace spanned by $e_1$ and $e_3$.
\end{example}

\begin{quote}\small
\textbf{Intuition.}
Spectral selection means projecting onto a discrete admissible spectrum $S$,
not necessarily one indexed by integers.
The set $S$ is chosen by invariance or admissibility constraints, not by fitting to data.
(The term ``quantization projector'' reflects physics usage where discrete spectra arise from quantization conditions.)
\end{quote}

\begin{remark}[Selection rule]
\lean{CoherenceCalculus.AppendixOperators.selectionRule}
In applications, $S$ is typically determined by conservation laws, symmetry constraints, or admissibility conditions (Section~\ref{sec:constraint-compiler}).
The spectral selector $Q_S(A)$ then projects onto the ``allowed'' eigenspaces.
\end{remark}

% ------------------------------------------------------------
\subsection{Coherence Trace (Horizon Trace and Supertrace)}
\label{subsec:coherence-trace}
% ------------------------------------------------------------

The coherence trace provides a canonical way to aggregate operator information across the observable subspace.

\begin{definition}[Horizon trace]
\label{def:horizon-trace}\lean{CoherenceCalculus.AppendixOperators.horizonTrace}
For a trace-class operator $A$ on $\cH$ (or any operator on a finite-dimensional $\cH$), define the \emph{horizon trace} at level $h$:
\[
\CCHTrace{h}{A} := \Tr(\CCRestr{A}{h}).
\]
\end{definition}

\begin{lemma}[Properties of horizon trace]
\label{lem:horizon-trace-properties}\lean{CoherenceCalculus.AppendixOperators.horizonTrace_add, CoherenceCalculus.AppendixOperators.horizonTrace_monotone_statement, CoherenceCalculus.AppendixOperators.horizonTrace_consistency}
The horizon trace satisfies:
\begin{enumerate}[label=(\roman*)]
\item \textbf{Linearity:} $\CCHTrace{h}{\alpha A + \beta B} = \alpha \CCHTrace{h}{A} + \beta \CCHTrace{h}{B}$.
\item \textbf{Monotonicity:} If $h \le h'$ (nested horizons), then for $A \succeq 0$, $\CCHTrace{h}{A} \le \CCHTrace{h'}{A}$.
\item \textbf{Consistency:} $\CCHTrace{h}{\CCPi{h}} = \dim(\CCHobs{h})$ (the dimension of the observable subspace).
\end{enumerate}
\end{lemma}

\begin{proof}
(i) Linearity of trace.
(ii) For $A \succeq 0$ and $h \le h'$, the range of $\CCPi{h}$ is contained in the range of $\CCPi{h'}$, so the sum defining $\CCHTrace{h}{A}$ is over a subset of the terms in $\CCHTrace{h'}{A}$.
(iii) $\Tr(\CCPi{h} \CCPi{h} \CCPi{h}) = \Tr(\CCPi{h}) = \dim(\CCHobs{h})$.
\end{proof}

\begin{definition}[Coherence supertrace]
\label{def:supertrace}\lean{CoherenceCalculus.AppendixOperators.coherenceSupertrace}
If $\cH = \bigoplus_{k=0}^{K} \cH^{(k)}$ is a graded Hilbert space with grading projections $P^{(k)}$, and $\CCPi{h}^{(k)} := \CCPi{h} P^{(k)}$ are the horizon projections restricted to grade $k$, define the \emph{coherence supertrace}:
\[
\mathrm{Str}_h(A) := \sum_{k=0}^{K} (-1)^k \Tr(\CCPi{h}^{(k)} A \CCPi{h}^{(k)}).
\]
Each term must be finite (finite-dimensional $\cH^{(k)}$ or trace-class restriction).
\end{definition}

\begin{lemma}[Supertrace linearity]
\label{lem:supertrace-linear}\lean{CoherenceCalculus.AppendixOperators.coherenceSupertrace_add}
The coherence supertrace is linear: $\mathrm{Str}_h(\alpha A + \beta B) = \alpha \, \mathrm{Str}_h(A) + \beta \, \mathrm{Str}_h(B)$.
\end{lemma}

\begin{proof}
Direct from linearity of trace on each graded component.
\end{proof}

\begin{quote}\small
\textbf{Intuition.}
The horizon trace $\CCHTrace{h}{A}$ replaces ad hoc ``sums over observable sectors'' with a canonical trace functional.
The supertrace extends this to graded/complex structures where alternating signs encode cohomological information.
\end{quote}

% ------------------------------------------------------------
\subsection{Aliasing Operator (Polar Decomposition of Leakage)}
\label{subsec:aliasing}
% ------------------------------------------------------------

The aliasing operator extracts the ``direction'' of leakage from the observable into the shadow subspace.

Recall the \emph{leakage operator} $\CCLeak{d}{h} := \CCLeakExpr{d}{h}$ (Definition~\ref{def:leakage-operator}), the component of $d$ that maps observable states into the shadow.

\begin{theorem}[Polar decomposition of leakage]
\label{thm:aliasing-polar}\lean{CoherenceCalculus.AppendixOperators.leakageMagnitudeSquared}
The leakage operator $B_h$ admits a polar decomposition
\[
B_h = U_h |B_h|,
\]
where:
\begin{itemize}[itemsep=0.2em]
\item $|B_h| := (B_h^* B_h)^{1/2}$ is the positive semidefinite ``magnitude'' of leakage.
\item $U_h$ is a partial isometry with initial space $\overline{\Img(|B_h|)}$ and final space $\overline{\Img(B_h)}$.
\end{itemize}
\end{theorem}

\begin{proof}
This is the standard polar decomposition theorem for bounded operators on Hilbert spaces.
\end{proof}

\begin{definition}[Aliasing operator]
\label{def:aliasing-operator}\lean{CoherenceCalculus.AppendixOperators.aliasingProjectorDomain}
The \emph{aliasing operator} is the partial isometry $U_h$ from the polar decomposition of $B_h$.
It encodes the ``canonical direction'' of leakage from observable to shadow.
\end{definition}

\begin{corollary}[Properties of aliasing]
\label{cor:aliasing-properties}\lean{CoherenceCalculus.AppendixOperators.leakageMagnitude_observable, CoherenceCalculus.AppendixOperators.isometry_property, CoherenceCalculus.AppendixOperators.leakage_norm_eq_abs_norm}
\begin{enumerate}[label=(\roman*)]
\item $B_h^* B_h = |B_h|^2 \succeq 0$.
\item $\|B_h x\|^2 = \langle x, B_h^* B_h \, x \rangle = \langle x, |B_h|^2 x \rangle$ for all $x \in \cH$.
\item $U_h$ preserves norms on the initial space $\Img(|B_h|)$: for all $z \in \Img(|B_h|)$, $\|U_h z\| = \|z\|$.
\item $U_h^* U_h$ is the projection onto $\overline{\Img(|B_h|)}$.
\end{enumerate}
\end{corollary}

\begin{proof}
Standard properties of polar decomposition.
\end{proof}

\begin{quote}\small
\textbf{Intuition.}
The aliasing operator $U_h$ is the ``direction'' of leakage; $|B_h|$ is its ``magnitude.''
Together they separate \emph{how much} leaks ($|B_h|$) from \emph{where} it goes ($U_h$).
\end{quote}

\begin{remark}[Connection to defect energy]
\lean{CoherenceCalculus.AppendixOperators.defect_energy_aliasing_direction}
The defect energy (Definition in Section~\ref{sec:hft}) satisfies
\[
\CCDefEnergy{h}{x} = \|\CCLeak{d}{h} x\|^2 = \||\CCLeak{d}{h}| x\|^2.
\]
The aliasing operator $U_h$ does not affect the magnitude of leakage, only its direction in the shadow subspace.
\end{remark}

\begin{example}[Finite-dimensional aliasing]
\lean{CoherenceCalculus.AppendixOperators.finite_quantization_example}
\label{ex:finite-aliasing}
Let $\cH = \bbR^2$, $\Pi = e_1 e_1^*$ (projection onto first coordinate), and $d = \begin{psmallmatrix} 0 & 1 \\ 1 & 0 \end{psmallmatrix}$.
Then $B = (I-\Pi)\, d\, \Pi = e_2 e_1^*$ (maps $e_1 \mapsto e_2$).
Here $|B| = e_1 e_1^*$ and $U = e_2 e_1^*$, so $U$ is the ``rotation'' from observable to shadow direction.
\end{example}
